\chapter{Prefácio}\label{preface}

\paragraph{Sobre este livro}

Como o título indica, este é um livro-texto sobre lógica formal. A lógica formal trata do estudo de um certo tipo de linguagem que, como qualquer linguagem, pode servir para expressar estados de coisas. É uma linguagem formal, isto é, suas expressões (como sentenças) são definidas formalmente. Isso a torna uma linguagem muito útil para tratar com precisão os estados de coisas que suas sentenças descrevem. Em particular, na lógica formal é impossível haver ambiguidade. O estudo dessas linguagens concentra-se na relação de consequência lógica entre sentenças, isto é, em quais sentenças decorrem de quais outras sentenças. A consequência lógica é central porque, ao compreendê-la melhor, podemos determinar quando certos estados de coisas devem ocorrer, desde que ocorram outros estados de coisas. Mas a consequência lógica não é a única noção importante. Também consideraremos a relação de satisfatibilidade, isto é, de não serem mutuamente contraditórias. Essas noções podem ser definidas semanticamente, usando definições precisas de consequência lógica baseadas em interpretações da linguagem — ou do ponto de vista da teoria da prova, usando sistemas formais de dedução.

A lógica formal é, naturalmente, uma subdisciplina central da filosofia, na qual é importante a relação lógica entre pressupostos e conclusões deles obtidas. Filósofos investigam as consequências de definições e pressupostos e avaliam essas definições e pressupostos com base em suas consequências. A lógica formal também é importante na matemática e na ciência da computação. Na matemática, as linguagens formais são usadas para descrever não estados de coisas da \textquotedblleft vida cotidiana\textquotedblright, mas estados de coisas matemáticos. Os matemáticos também se interessam pelas consequências de definições e pressupostos e, para eles, é igualmente importante estabelecer essas consequências (que chamam de \textquotedblleft teoremas\textquotedblright) por meio de métodos completamente precisos e rigorosos. A lógica formal fornece tais métodos. Na ciência da computação, a lógica formal é aplicada para descrever o estado e os comportamentos de sistemas computacionais, como circuitos, programas, bancos de dados etc. Métodos de lógica formal também podem ser usados para estabelecer consequências dessas descrições, como verificar se um circuito não contém erros, se um programa faz aquilo que deve fazer, se um banco de dados é consistente ou se algo é verdadeiro a respeito dos dados nele contidos.

O livro está dividido em nove partes. \Cref{ch.intro} introduz o tema e as noções de lógica de maneira informal, sem ainda introduzir uma linguagem formal. \Crefrange{ch.TFL}{ch.NDTFL} tratam de linguagens verofuncionais. Nelas, as sentenças são formadas a partir de sentenças básicas usando diversos conectivos (\textquotedblleft ou\textquotedblright, \textquotedblleft e\textquotedblright, \textquotedblleft não\textquotedblright, \textquotedblleft se \dots então\textquotedblright) que simplesmente combinam sentenças em outras mais complexas. Discutimos noções lógicas como a consequência lógica de duas maneiras: semanticamente, usando o método das tabelas-verdade (em \cref{ch.TruthTables}), e do ponto de vista da teoria da prova, usando um sistema de derivações formais (em \cref{ch.NDTFL}). \Crefrange{ch.FOL}{ch.NDFOL} tratam de uma linguagem mais complexa: a lógica de primeira ordem. Além dos conectivos da lógica verofuncional, ela inclui também nomes, predicados, identidade e os chamados quantificadores. Esses elementos adicionais tornam essa linguagem muito mais expressiva que a linguagem verofuncional, e passaremos um bom tempo investigando precisamente quanto é possível expressar nela. Novamente, as noções lógicas da linguagem de primeira ordem são definidas semanticamente, usando interpretações, e do ponto de vista da teoria da prova, usando uma versão mais complexa do sistema de derivações formais introduzido em \cref{ch.NDTFL}. \Cref{ch.ML} discute a extensão da LVF por operadores não verofuncionais de possibilidade e necessidade: a lógica modal. \Cref{ch.normalform} aborda dois tópicos avançados: formas normais conjuntivas e disjuntivas e a completude funcional dos conectivos verofuncionais, além da correção da dedução natural para a LVF.

Nos apêndices, você encontrará uma discussão sobre notações alternativas para as linguagens tratadas neste texto, sistemas alternativos de derivação e uma referência rápida que lista a maioria das regras e definições importantes. Os termos centrais estão listados em um glossário ao final do livro.

\paragraph{Créditos} Este livro é baseado em um texto originalmente escrito por
P.~D. Magnus, na versão revisada e ampliada por Tim Button. Ele
também inclui algum material (principalmente exercícios) de J.~Robert Loftis.
O material em \cref{ch.ML} baseia-se em notas de Robert Trueman (mas foi
reescrito para usar as regras originais de dedução natural de Fitch para a
lógica modal), e o material em
\cref{c:NormalForms,c:FunctionalCompleteness,ch:Soundness} baseia-se em dois
capítulos do texto aberto \textit{Metatheory}, de Tim Button. Aaron
Thomas-Bolduc e Richard Zach combinaram elementos desses textos
na presente versão, alteraram parte da terminologia e dos
exemplos, reescreveram algumas seções e acrescentaram material próprio. Em
particular, Richard Zach reescreveu \cref{s:Arguments,s:Valid} e acrescentou
\cref{s:AmbiguityTFL,s:stratTFL,s:ambiguityFOL,ch:PropRelations,ch:equivalences}.
A partir da edição do outono de 2019, a parte sobre LPO usa a sintaxe mais
comum em textos avançados (como os baseados no
\href{https://openlogicproject.org/}{Open Logic Project}), nos quais os
argumentos dos símbolos de predicado são colocados entre parênteses (isto é,
\textquotedblleft $R(a,b)$\textquotedblright{} em vez de \textquotedblleft $Rab$\textquotedblright{}). Esta versão de \forallx{} também usa a
definição padrão de sintaxe para LPO, que permite a quantificação vazia, e as
regras originais de introdução e eliminação da dedução natural de Gentzen (isto
é, as regras de eliminação da dupla negação e do terceiro excluído são regras
derivadas, e a prova indireta é a única regra clássica). O texto resultante é
licenciado sob uma licença Creative Commons Atribuição 4.0. Existem
\href{https://github.com/OpenLogicProject/OpenLogic/wiki/Other-Logic-Textbooks}{vários
outros} \textquotedblleft remixes\textquotedblright{} de \forallx, incluindo traduções desta versão.

\paragraph{Notas para docentes} O material deste livro é adequado para uma introdução à lógica formal ao longo de um semestre. Eu cubro \crefrange{ch.intro}{ch.NDFOL}, além de \cref{c:NormalForms,c:FunctionalCompleteness,ch:equivalences}, em 12 semanas, embora deixe de fora as tabelas-verdade parciais e as regras de inferência derivadas.

A versão mais recente deste livro está disponível em PDF em
\href{https://forallx.openlogicproject.org}{forallx.openlogicproject.org},
mas muda com frequência. A licença CC BY dá a você o direito de
baixar e distribuir o livro por conta própria. Para garantir que todos os
seus estudantes tenham a mesma versão do livro durante o período em que você
o estiver utilizando, faça isso: envie o PDF escolhido para o seu AVA, em vez
de simplesmente fornecer aos estudantes o link. O livro também está disponível em
uma
\href{https://forallx.openlogicproject.org/html/}{versão HTML}. Essa
versão foi preparada com
\href{https://forallx.openlogicproject.org/html/A4.html}{atenção especial às
questões de acessibilidade} e deve ser mais adequada para pessoas que dependem
de leitores de tela, por exemplo, estudantes com baixa visão ou perda total da
visão. Você pode baixar um pacote SCORM que inclui as versões HTML e PDF para
enviá-las ao seu AVA. Você também pode pedir que os PDFs sejam impressos por sua
livraria, mas algumas livrarias talvez consigam comprar e manter em estoque os
livros de capa flexível disponíveis na Amazon.

Observe que as soluções de muitos exercícios do livro também estão disponíveis
no site acima (para todos, inclusive seus estudantes). As soluções também
fazem parte da versão HTML e do pacote SCORM.

A sintaxe e os sistemas de prova (exceto os de lógica modal) são apoiados pelo
aplicativo gratuito de ensino de lógica on-line \textit{Carnap}, de Graham
Leach-Krouse (\href{https://carnap.io}{carnap.io}). Ele permite o envio e a
correção automática de exercícios de simbolização, tabelas-verdade e provas de
dedução natural. Docentes em carnap.io poderão encontrar
\href{https://carnap.io/shared/rzach@ucalgary.ca/forall%20x:%20Calgary.md}{amostras de exercícios adicionais}
que talvez queiram adaptar ou atribuir tal como estão.
